\subsection{Nonconjugate block calculations: what each approximation establishes}
\label{sec:nonconjugate-derivations}
For a canonical Binomial--logistic block with a Gaussian prior on its intercept, the exact local evidence is a one-dimensional integral. It provides a numerical reference against which deterministic approximations can be tested before the resulting segment scores enter the same partition recursion. A local approximation formula, its variational lower bound, and a rigorous numerical error bound are three different objects.

\begin{lemma}[Finite, unique logistic-block mode under a proper normal prior]
\label{lem:logistic-mode-paper}
Let $y_r\mid\theta\sim\operatorname{Binom}(m_r,\sigma(\theta))$ independently, with $m_r\ge0$, $\sigma(\theta)=(1+e^{-\theta})^{-1}$, and $\theta\sim\mathcal N(\nu,\rho^2)$ for finite $\rho^2>0$. Then the block log posterior is strictly concave and has a unique finite maximizer, including blocks containing only successes or only failures.
\end{lemma}
\begin{proof}
\textbf{Step 1 (posterior objective).} Apart from its Binomial coefficients, the log posterior is
\[
\Psi_B(\theta)=Y_B\theta-M_B\log(1+e^\theta)
-\frac{(\theta-\nu)^2}{2\rho^2}+\mathrm{constant},
\]
where $Y_B=\sum y_r$ and $M_B=\sum m_r$.
\textbf{Step 2 (strict curvature).} Direct differentiation gives
\[
\Psi_B''(\theta)=-M_B\sigma(\theta)\{1-\sigma(\theta)\}-\rho^{-2}<0
\]
for every finite $\theta$, even when $M_B=0$.
\textbf{Step 3 (existence, not just uniqueness).} The Binomial log-likelihood is bounded above by the sum of its log combinatorial coefficients. The negative quadratic prior term tends to $-\infty$ as $|\theta|\to\infty$. Hence $\Psi_B$ is coercive, attains a finite maximum by continuity, and strict concavity makes that maximizer unique. This explicit tail condition is why the conclusion does not follow from strict concavity alone under an arbitrary prior.
\end{proof}

\subsubsection{Laplace evaluation and its limitations}
Let $\widehat\theta_B$ be the mode and $H_B=-\Psi_B''(\widehat\theta_B)>0$. The first-order Laplace approximation is
\begin{equation}
\log\widehat A_B^{\rm Lap}=\Psi_B(\widehat\theta_B)
+\tfrac12\log(2\pi)-\tfrac12\log H_B.
\label{eq:laplace-paper}
\end{equation}
This is a local second-order approximation to the integral, \emph{not} a nonasymptotic upper bound on the absolute log-evidence error. A controlled asymptotic remainder would require derivatives beyond second order, concentration of the integral near the mode, and bounds on the tails outside that neighborhood. Numerical evaluation should use safeguarded Newton steps, check finite curvature and convergence, and fall back to adaptive one-dimensional integration on failure. The software's local curvature diagnostic does not by itself establish Assumption~\ref{ass:block-error}.

\subsubsection{Jaakkola--Jordan lower bound}
\begin{proposition}[Valid logistic tangent bound and block-evidence lower bound]
\label{prop:jj-paper}
For $\xi>0$ put $\lambda(\xi)=\tanh(\xi/2)/(4\xi)$ and $\sigma(x)=(1+e^{-x})^{-1}$. Then, for every real $x$,
\begin{equation}
\log\sigma(x)\ge\log\sigma(\xi)+\frac{x-\xi}{2}
-\lambda(\xi)(x^2-\xi^2).
\label{eq:jj-correct-paper}
\end{equation}
For a $\operatorname{Binom}(m_r,\sigma(\theta))$ observation, this quadratic lower bound has linear coefficient $y_r-m_r/2$ and quadratic coefficient $-m_r\lambda(\xi_r)$. Combining these terms with a Gaussian prior and integrating the resulting Gaussian kernel gives a deterministic \emph{lower} bound on $\log A^{(0)}_B$; this bound alone supplies no upper bound on the unknown error gap.
\end{proposition}
\begin{proof}
\textbf{Step 1 (rewrite the log sigmoid).} The identity $\log\sigma(x)=x/2-\log2-\log\cosh(x/2)$ separates a linear term from an even function of $x$.
\textbf{Step 2 (bound the even term).} Define $f(v)=\log\cosh(\sqrt v/2)$ for $v\ge0$. Its derivative for $v>0$ is $f'(v)=\tanh(\sqrt v/2)/(4\sqrt v)$, which decreases with $v$. Thus $f$ is concave and lies below its tangent at $v=\xi^2$: $f(x^2)\le f(\xi^2)+\lambda(\xi)(x^2-\xi^2)$.
\textbf{Step 3 (recover the inequality).} Substitute this tangent inequality into the expression for $\log\sigma(x)$, then use the identity at $x=\xi$; Equation~\eqref{eq:jj-correct-paper} follows, with equality at $x=\pm\xi$ in the even part.
\textbf{Step 4 (apply to Binomial data).} Write $\log p(y_r\mid\theta)=\log\binom{m_r}{y_r}+(y_r-m_r/2)\theta-m_r\log[2\cosh(\theta/2)]$. Apply the tangent bound to its even part and sum over observations. Exponentiating and integrating an everywhere-lower-bounded nonnegative integrand preserves the ordering. Thus $\log\widehat A_B^{\rm JJ}\le\log A_B$, but the value of $\log A_B-\log\widehat A_B^{\rm JJ}$ has not been upper-bounded by this argument.
\end{proof}

\subsubsection{P\'olya--Gamma augmentation and mean-field computation}
\begin{proposition}[Augmentation identity and Gaussian coordinate updates]
\label{prop:pg-paper}
For a Binomial likelihood without its $\theta$-independent coefficient, let $b_r=m_r$, $\kappa_r=y_r-m_r/2$ and denote the P\'olya--Gamma density by $p_{\rm PG}(\omega\mid b,0)$. Then
\begin{equation}
\frac{e^{y_r\theta}}{(1+e^\theta)^{m_r}}
=2^{-b_r}e^{\kappa_r\theta}
\int_0^\infty e^{-\omega\theta^2/2}p_{\rm PG}(\omega\mid b_r,0)\dd\omega.
\label{eq:pg-identity-paper}
\end{equation}
With $\theta\sim\mathcal N(\nu,\rho^2)$, the mean-field family $q(\theta)\prod_rq(\omega_r)$ admits coordinate updates
\begin{align}
q(\theta)&=\mathcal N(\mu_q,V_q),
& V_q^{-1}&=\rho^{-2}+\sum_r\mathbb E_q[\omega_r],\nonumber\\
\mu_q&=V_q\left(\nu/\rho^2+\sum_r\kappa_r\right),
& q(\omega_r)&={\rm PG}(b_r,c_r),
& c_r^2&=\mu_q^2+V_q,
\label{eq:pg-updates-paper}
\end{align}
where $\mathbb E_q[\omega_r]=b_r\tanh(c_r/2)/(2c_r)$ with its continuous limit at $c_r=0$. If an additional \emph{whole-likelihood} power $a_r$ is used, the identity requires $(b_r,\kappa_r)=(a_rm_r,a_r(y_r-m_r/2))$; multiplying an independent ${\rm PG}(m_r,0)$ variable by $a_r$ is not equivalent.
\end{proposition}
\begin{proof}
\textbf{Step 1 (factor logistic likelihood).} The identity $1+e^\theta=2e^{\theta/2}\cosh(\theta/2)$ converts the Binomial likelihood into $2^{-m_r}e^{\kappa_r\theta}\cosh(\theta/2)^{-m_r}$.
\textbf{Step 2 (introduce the PG mixture).} The P\'olya--Gamma Laplace-transform identity is $\cosh(\theta/2)^{-b}=\int e^{-\omega\theta^2/2}p_{\rm PG}(\omega\mid b,0)\dd\omega$ for $b>0$. Substitution proves Equation~\eqref{eq:pg-identity-paper} (with the $b=0$ contribution interpreted as one).
\textbf{Step 3 (update the Gaussian factor).} The augmented log joint is quadratic in $\theta$ conditional on the PG variables. Taking the expectation over $q(\omega)$ and adding the Gaussian prior gives precision $\rho^{-2}+\sum_r\mathbb E_q[\omega_r]$ and linear coefficient $\nu/\rho^2+\sum_r\kappa_r$, proving the Gaussian update.
\textbf{Step 4 (update the latent factors).} The coordinate log density for $q(\omega_r)$ is $\log p_{\rm PG}(\omega_r\mid b_r,0)-\omega_r\mathbb E_q[\theta^2]/2$ plus a constant. Exponential tilting gives ${\rm PG}(b_r,c_r)$, $c_r^2=\mathbb E_q[\theta^2]$, and the standard PG mean gives Equation~\eqref{eq:pg-updates-paper}. Each exact coordinate maximization cannot lower its \emph{properly evaluated} ELBO; it does not certify the gap to the exact evidence.
\textbf{Step 5 (external powers).} Raising the complete Binomial factor to $a_r$ changes both the hyperbolic-cosine exponent and linear coefficient, hence the PG shape and $\kappa_r$ as stated. A scaled PG random variable does not have this tilted law.
\end{proof}

For the intercept-only logistic--Gaussian block, the optimized Jaakkola--Jordan quadratic bound and the corresponding PG mean-field Gaussian updates can be algebraically equivalent. For the present numerical comparison, the \texttt{jj} and \texttt{pg\_vb} labels correspond to the same implementation and hence represent a single computational method. A distinct P\'olya--Gamma implementation would require separate numerical validation.

\subsubsection{Expectation propagation and numerical reference}
Expectation propagation approximates likelihood contributions with Gaussian sites and iterates cavity/tilted-distribution moment matching. A valid approximation to the \emph{segment marginal likelihood} must include the global Gaussian normalizer and the site/tilted normalizing corrections, not merely integrate a normalized approximate posterior (whose integral is one). The EP estimates are not used to establish integrated-likelihood accuracy in the absence of a separately validated normalizer, damping scheme, convergence criterion, and reference comparison. Adaptive one-dimensional quadrature, with a documented integration tolerance and tail or endpoint control, remains the numerical reference for low-dimensional blocks.
